\documentclass[a4paper,12pt]{article}
\usepackage{../packages/coursCollege}
\newcommand{\Chapitre}{Limites}
\renewcommand{\path}{../../}

\renewcommand{\cours}{3MA1~--~EG~--~ns~--~2025-2026}
\begin{document}
\tocloftpagestyle{fancy}
% Reduce space between section entries
\setlength{\cftbeforesecskip}{2pt}

% Reduce indentation for section entries
\setlength{\cftsecindent}{1em}

\begin{center}
	{\bfseries \Huge Chapitre 3~: Fonctions trigonométriques}

\bigskip

\tableofcontents
\newpage
\section{Limites trigonométriques}
Le théorème suivant est important et sa démonstration sera faite en activité. 
\medskip
\begin{thm}[label=thm:sinxx]
	\smallskip
$\displaystyle\lim_{x \to 0} \dfrac{\sin(x)}{x}$
	\tcblower
On a  \[\displaystyle\lim_{x \to 0} \dfrac{\sin(x)}{x} = 1\]
\end{thm}

\begin{methode}
\tcblower
Pour lever une indétermination trigonométrique, on se ramène au cas du théorème \ref{thm:sinxx} afin d'utiliser la formule $\displaystyle\lim_{\star \to 0} \dfrac{\sin(\star)}{\star} = 1$ où $\star$ peut être une expression quelconque.
\end{methode}

\begin{exemple}
	$\displaystyle\lim_{x \to 0} \dfrac{\sin(-20x)}{x}$
	\tcblower
Calculer $\displaystyle\lim_{x \to 0} \dfrac{\sin(-20x)}{x}$
\smallskip
\begin{explanation}
	$\displaystyle\lim_{x \to 0} \dfrac{\sin(-20x)}{x} = $ « $\dfrac{0}{0}$ » & on substitue $x$ par $0$ et on constate un « type $\dfrac{0}{0}$ »\\
	$\displaystyle\lim_{x \to 0} \dfrac{\sin(-20x)}{x}$ & on observe que l'expression qui « pilote » le calcul est $-20x$\\
	$= \displaystyle\lim_{-20x \to 0} \dfrac{-20\sin(-20x)}{-20x}$& on fait apparaître cette expression au dénominateur et comme le terme de la limite\\
	$= -20 \displaystyle\lim_{-20x \to 0} \dfrac{\sin(-20x)}{-20x}$& on utilise le théorème \ref{thm:sinxx}\\
	$= -20 \cdot 1$&on utilise $\displaystyle\lim_{\star \to 0} \dfrac{\sin(\star)}{\star} = 1$\\
$= -20$\\
\end{explanation}
\end{exemple}
\begin{remarque}
	\tcblower
On est parfois amené à devoir utiliser des propriétés trigonométriques ou des idées plus subtiles ...
\end{remarque}
\begin{exemple}
$\displaystyle\lim_{x \to 0} \dfrac{1-\cos(x)}{3x}$
	\tcblower
Calculer $\displaystyle\lim_{x \to 0} \dfrac{1-\cos(x)}{3x}$
\begin{explanation}[0.6]
	$\displaystyle\lim_{x \to 0} \dfrac{1-\cos(x)}{3x}$& on substitue $x$ par $0$ et on constate un « type $\dfrac{0}{0}$ »\\
	$= \displaystyle\lim_{x \to 0} \dfrac{1-\cos(x)}{3x^2} \cdot \dfrac{1+\cos(x)}{1+\cos(x)}$& on multiplie par le conjugué\\
	$= \displaystyle\lim_{x \to 0} \dfrac{(1-\cos(x))(1+\cos(x))}{3x^2(1+\cos(x))}$& on effectue la multiplication des fractions\\
	$= \displaystyle\lim_{x \to 0} \dfrac{1-\cos^2(x)}{3x^2(1+\cos(x))}$& on utilise la 3e identité remarquable\\
	$= \displaystyle\lim_{x \to 0} \dfrac{\sin^2(x)}{3x^2(1+\cos(x))}$& on utilise la propriété trigonométrique $\sin^2(x)+\cos^2(x)=1$\\
	$= \displaystyle\lim_{x \to 0} \dfrac{1}{3} \cdot \dfrac{\sin^2(x)}{x^2} \cdot \dfrac{1}{1+\cos(x)}$& on sépare en produit de fractions\\
	$= \dfrac{1}{3} \displaystyle\lim_{x \to 0} \dfrac{\sin(x)}{x} \displaystyle\lim_{x \to 0} \dfrac{\sin(x)}{x} \displaystyle\lim_{x \to 0} \dfrac{1}{1+\cos(x)}$& on utilise les propriétés des limites, ce qu'on peut faire si les deux limites existent\\
	$= \dfrac{1}{3} \cdot 1 \cdot 1 \cdot \dfrac{1}{1+\cos(0)}$& on utilise $\displaystyle\lim_{\star \to 0} \dfrac{\sin(\star)}{\star} = 1$ et on calcule la 3e limite facilement\\
$= \dfrac{1}{3} \cdot \dfrac{1}{2} = \dfrac{1}{6}$
\end{explanation}
\end{exemple}

\section{Dérivées des fonctions trigonométriques}
Théorème [Dérivée du sinus]
\begin{thm}
Dérivées des fonctions trigonométriques
\tcblower
On a 
\begin{itemize}
	\item $[\sin(x)]' = \cos(x)$
	\item $[\cos(x)]' = -\sin(x)$
	\item $[\tan(x)]' = 1 + \tan^2(x) = \dfrac{1}{\cos^2(x)}$
\end{itemize}
\end{thm}


\begin{notation}
	\tcblower
De même que pour $[\sin(x)]^2 = \sin^2(x)$ on note plus simplement $[\sin(x)]' = \sin'(x)$, $[\cos(x)]' = \cos'(x)$ et $[\tan(x)]' = \tan'(x)$.
\end{notation}
\begin{exemple}
	$[\sin(-20x^3 + x)]'$
	\tcblower
On utilise la formule de dérivation de la composition de fonctions. 

$\sin(-20x^3 + x) = g(f(x))$ avec $f(x) = -20x^3 + x$ et $g(x) = \sin(x)$

$[\sin(-20x^3 + x)]' = \cos(-20x^3 + x) \cdot (-20x^3 + x)' = \cos(-20x^3 + x) \cdot (-60x^2 + 1)$


\end{exemple}

\begin{exemple}
	$[\cos^5(x)]'$
	\tcblower
On utilise la formule de dérivation de la composition de fonctions. 

$\cos^5(x) = [\cos(x)]^5 = g(f(x))$ avec $f(x) = \cos(x)$ et $g(x) = x^5$

$[\cos^5(x)]' =$
$[[\cos(x)]^5]' = 5[\cos(x)]^4 \cdot [\cos(x)]' = 5[\cos(x)]^4 \cdot [-\sin(x)] = 5[-\cos(x)]^4 \sin(x)$
\end{exemple}

\begin{exemple}
$[\sin^3(x^2 - 1)]'$
	\tcblower
On utilise la formule de dérivation de la composition de fonctions. 

$\sin^3(x^2 - 1) = [\sin(x^2 - 1)]^3 = g(f(x))$ avec $f(x) = \sin(x^2 - 1)$ et $g(x) = x^5$

$[[\sin(x^2 - 1)]^3]' = 3[\sin(x^2 - 1)]^2 \cdot [\sin(x^2 - 1)]'$

et on doit recommencer pour dériver $[\sin(x^2 - 1)]'$ :

$[\sin(x^2 - 1)]' = \cos(x^2 - 1) \cdot (x^2 - 1)' = \cos(x^2 - 1) \cdot (2x)$

d'où enfin : $[[\sin(x^2 - 1)]^3]' = 3[\sin(x^2 - 1)]^2 \cos(x^2 - 1) \cdot 2x = 6x\sin^2(x^2 - 1)\cos(x^2 - 1)$
\end{exemple}





Librement adapté de ...

{\bfseries Source principale : Sesamath 3e année, 2023--2024}

\end{document}

